\section{Numerics}
\label{sec:numerics}

Every number in this section was computed for this paper: the exact quantities by evaluating the formulas, the Monte Carlo quantities by simulating the network with a fixed seed and comparing with the formulas they are supposed to obey. Where a quantity was sampled rather than enumerated the table says so, since a sampled supremum is a lower bound.

\subsection{The orthant count}

\Cref{tab:cover} draws a random subspace $U$ of dimension $d$ in $\R^M$ and counts the open orthants it meets. In the plane the count is exact, by enumerating the arcs between the angles at which the coordinate functionals vanish; for $d\ge3$ the count is a lower bound from sampled directions. The bound of \cref{lem:cover} is attained whenever the count is exact, in agreement with the general-position remark.

\begin{table}[h]\centering
\begin{tabular}{rrrrrl}
\toprule
$M$ & $d$ & $2^M$ & $2\sum_{j<d}\binom{M-1}j$ & observed $r(U)$, mean (max) & method\\
\midrule
6 & 2 & 64 & 12 & 12.0 (12) & exact\\
8 & 2 & 256 & 16 & 16.0 (16) & exact\\
10 & 2 & 1024 & 20 & 20.0 (20) & exact\\
20 & 2 & 1048576 & 40 & 40.0 (40) & exact\\
6 & 3 & 64 & 32 & 31.9 (32) & sampled\\
9 & 3 & 512 & 74 & 73.4 (74) & sampled\\
12 & 3 & 4096 & 134 & 131.8 (134) & sampled\\
8 & 4 & 256 & 128 & 120.8 (127) & sampled\\
\bottomrule
\end{tabular}
\caption{Orthants met by a random $d$-dimensional subspace of $\R^M$ against the bound of \cref{lem:cover}; 30 subspaces per row for $d=2$, 12 for $d\ge3$.}
\label{tab:cover}
\end{table}

For the network itself, the identity \cref{eq:orbit} with $F=1$, summed over all masks, says that the expected number of nonempty strict regions is $\sum_D\E\,r(\im A_D)/2^{nL}\le C_{n,L}$. \Cref{tab:regions} compares the number of realized masks with $C_{n,L}$ and with the number $2^{nL}$ of formal masks. Even at width two the realized count stays far below the formal count, and below $C_{n,L}$.

\begin{table}[h]\centering
\begin{tabular}{rrrrrl}
\toprule
$n$ & $L$ & formal masks $2^{nL}$ & $C_{n,L}$ & realized masks, mean (min--max) & method\\
\midrule
2 & 1 & 4 & 4 & 4.0 (4--4) & exact\\
2 & 2 & 16 & 8 & 4.9 (4--6) & exact\\
2 & 3 & 64 & 12 & 5.0 (4--7) & exact\\
2 & 4 & 256 & 16 & 5.5 (4--8) & exact\\
2 & 5 & 1024 & 20 & 5.3 (4--8) & exact\\
2 & 6 & 4096 & 24 & 5.5 (4--10) & exact\\
3 & 2 & 64 & 32 & 15.5 (8--20) & sampled\\
3 & 3 & 512 & 74 & 20.3 (11--28) & sampled\\
4 & 2 & 256 & 128 & 47.1 (24--71) & sampled\\
\bottomrule
\end{tabular}
\caption{Nonempty strict regions of the random network (60 draws per row at $n=2$, 25 otherwise).}
\label{tab:regions}
\end{table}

At fixed $L$ the per-width exponent of $C_{n,L}$ converges to $LH(1/L)$ (\cref{lem:C-rate}). At $n=40$: for $L=3$, $\frac1n\log C_{n,L}=1.851$ against $LH(1/L)=1.910$ and $L\log2=2.079$; for $L=10$, $3.146$ against $3.251$ and $6.931$; for $L=100$, $5.434$ against $5.600$ and $69.31$. The entropy bound of \cref{lem:C-entropy} is within a few percent of the exact count per unit width already at width 40, and the saving against the $2^{nL}$ formal masks is the whole difference between $\log L+1$ and $L\log2$.

\subsection{The realizable-mask identity}

\Cref{tab:orbit} is the paired test of \cref{prop:orbit} at width $n=2$ described in \cref{sec:gauge}: for a fixed mask $D$, the same weight draws are used on both sides, $r(\im A_D)$ is computed exactly by arc enumeration, and the paired difference is compared with its standard error. Two functionals are used, $F\equiv1$ and $F=\min(\op{J_D},1)$, both bounded, so the standard error is meaningful.

\begin{table}[h]\centering
\begin{tabular}{rlrrrrr}
\toprule
$L$ & $F$ & draws & $\E[F\one_{\{R_D\ne\varnothing\}}]$ & $\E[F\,r(\im A_D)/2^{nL}]$ & ratio & paired $z$\\
\midrule
2 & $1$ & 200000 & 0.24878 & 0.25000 & 0.995 & $-1.27$\\
3 & $1$ & 200000 & 0.06283 & 0.06250 & 1.005 & $+0.59$\\
4 & $1$ & 150000 & 0.01589 & 0.01563 & 1.017 & $+0.83$\\
5 & $1$ & 120000 & 0.00390 & 0.00391 & 0.998 & $-0.04$\\
3 & $\min(\op{J_D},1)$ & 200000 & 0.02307 & 0.02296 & 1.005 & $+0.41$\\
4 & $\min(\op{J_D},1)$ & 150000 & 0.00492 & 0.00493 & 0.999 & $-0.03$\\
\bottomrule
\end{tabular}
\caption{The exact orbit identity \cref{eq:orbit} at width 2, the mask $D$ fixed, the subspace $\im A_D$ enumerated exactly.}
\label{tab:orbit}
\end{table}

Two more checks that the proof leans on directly were run alongside: summing the indicator $\one_{\{S_\Sigma U\cap O_D\ne\varnothing\}}$ over the whole gauge group returned exactly $r(U)$ in every draw, and $\op{J_D}$ was unchanged by a random gauge in every draw, to machine precision.

The all-active formula \cref{eq:all-active} gives $\PP(R_D\ne\varnothing)=\tfrac12$ at $(n,L)=(2,2)$ and $(3,2)$, $0.1445$ at $(3,3)$, $0.0898$ at $(5,3)$, $0.0307$ at $(10,3)$, $0.00432$ at $(20,3)$ and $3.0\times10^{-18}$ at $(10,10)$: the chance that a deep random network has an input at which every unit fires at every layer is exactly computable and small.

\subsection{The moment formulas}

\Cref{tab:radial} checks the exact radial factor \cref{eq:radial} and the mixture identity \cref{eq:mixture} by brute force.

\begin{table}[h]\centering\small
\begin{tabular}{llrrr}
\toprule
identity & parameters & Monte Carlo & exact & ratio\\
\midrule
\cref{eq:radial} & $n=4$, $L=3$, $p=1$, $\alpha=1.3$, masks $1101,1001,1110$ & 0.6247 & 0.6179 & 1.011\\
\cref{eq:radial} & $n=4$, $L=3$, $p=2$, same masks & 2.063 & 2.121 & 0.973\\
\cref{eq:radial} & $n=6$, $L=2$, $p=1.5$, $\alpha=1.3$, masks $101101,011011$ & 0.9022 & 0.8987 & 1.004\\
\cref{eq:mixture} & $n=6$, $p=1$, $\alpha=1.3$ & 41.65 & 41.60 & 1.001\\
\cref{eq:mixture} & $n=6$, $p=2.5$, $\alpha=1.3$ & 61.99 & 61.93 & 1.001\\
\cref{eq:mixture} & $n=12$, $p=1.7$, $\alpha=2$ & 5087 & 5079 & 1.002\\
\bottomrule
\end{tabular}
\caption{The fixed-vector radial moments and the scalar mixture identity, $2\times10^5$ draws each; the masks are the diagonals of $D_1,\dots,D_L$; for \cref{eq:mixture} both sides are multiplied by $2^n$.}
\label{tab:radial}
\end{table}

The master inequality \cref{eq:master} is far from tight at small width, and the table says by how much: at $(n,L,\alpha,p)=(2,2,1,1)$ the left side $\E K^{2p}$ is $1.51$ against a bound of $200$; at $(2,3,1,1.5)$, $1.94$ against $557$; at $(2,2,0.6,2)$, $0.88$ against $318$. The slack is the product of the net factor $5^n4^p$ and of $C_{n,L}$ counting every orthant a subspace could meet rather than the ones it does. The inequality is designed for the exponential scale, where these factors cost a fixed rate per unit width, and it is there that it is used.

The $\tfrac12$-net of \cref{lem:net} is also loose by a fixed exponential factor: a greedy maximal $\tfrac12$-separated set on $S^{n-1}$ had $2,9,34,102,279,676$ points for $n=1,\dots,6$ against the bound $5^n$; the ratio behaves like $e^{-0.5n}$. Neither slack affects the threshold, which is decided by the $L$-fold radial factor.

\subsection{The finite-width bound}

\Cref{tab:threshold} lists the threshold $\alpha^*_L$ of \cref{eq:alpha-star}, the largest $\alpha$ at which the exact budget $B_L$ is negative for some $\theta$, together with the asymptotic formula of \cref{thm:main}(iv) evaluated at the same $L$. The asymptotic formula overstates the bound at moderate depth: at $L=100$ it reads $1.01$ with the $o(1)$ dropped, or $0.81$ with the constant $A=1+\log5+1$ kept, while the budget delivers $0.68$. At $L=10^5$ the three agree to a few parts in a thousand.

\begin{table}[h]\centering
\begin{tabular}{rrrr}
\toprule
$L$ & $\alpha^*_L$ (exact budget) & $2e^{-\sqrt{10\log L/L}}$ & $2e^{-\sqrt{10(\log L+A)/L}}$, $A=2+\log5$\\
\midrule
10 & 0.0412 & 0.4386 & 0.1758\\
30 & 0.2502 & 0.6896 & 0.4336\\
100 & 0.6778 & 1.0146 & 0.8080\\
300 & 1.0822 & 1.2932 & 1.1457\\
1000 & 1.4261 & 1.5378 & 1.4461\\
3000 & 1.6390 & 1.6986 & 1.6428\\
$10^4$ & 1.7875 & 1.8170 & 1.7859\\
$10^5$ & 1.9254 & 1.9333 & 1.9237\\
$10^6$ & 1.9744 & 1.9766 & 1.9738\\
\bottomrule
\end{tabular}
\caption{The threshold the finite theorem delivers at depth $L$, against the large-$L$ formula.}
\label{tab:threshold}
\end{table}

The other reading of the same budget is the depth $L_0(\alpha)$ at which it first turns negative: $237$ for $\alpha=1$, $897$ for $1.4$, $4590$ for $1.7$, $52\,841$ for $1.9$, $237\,405$ for $1.95$ and $7.35\times10^6$ for $1.99$. The convergence $\alpha_L\to2$ is proved, and it is slow.

\Cref{tab:rate} gives the width rate $c_L(\alpha)$ of \cref{cor:width} and, in parentheses, the width at which the bound $2e^{-c_L(\alpha)n}$ drops below $0.01$. Once the depth is past $L_0(\alpha)$ the rate grows linearly in $L$, and a width of one or two already makes the bound tiny; at $\alpha=1$, $L=300$ the bound is $5.6\times10^{-9}$ at $n=10$ and $3.4\times10^{-43}$ at $n=50$, at $L=1000$ it is $3.8\times10^{-109}$ at $n=10$. The finite-width statement is therefore not weak once it applies; what limits it is the depth $L_0(\alpha)$, not the width.

\begin{table}[h]\centering
\begin{tabular}{rrrrrrr}
\toprule
$\alpha\backslash L$ & 100 & 300 & 1000 & 3000 & $10^4$ & $10^5$\\
\midrule
0.5 & 3.63 (2) & 24.5 (1) & 100.4 (1) & 319.5 (1) & 1089 (1) & 10996 (1)\\
1.0 & -- & 1.97 (3) & 25.0 (1) & 93.3 (1) & 334.8 (1) & 3453 (1)\\
1.5 & -- & -- & -- & 10.6 (1) & 59.2 (1) & 696.5 (1)\\
1.8 & -- & -- & -- & -- & -- & 90.5 (1)\\
1.9 & -- & -- & -- & -- & -- & 11.4 (1)\\
\bottomrule
\end{tabular}
\caption{The width rate $c_L(\alpha)$, with the width at which $2e^{-c_L(\alpha)n}\le0.01$ in parentheses; a dash means the budget is not negative at that depth.}
\label{tab:rate}
\end{table}

On the other side, \cref{thm:above2} gives $\PP(K_{n,L}(\alpha)\le1)\le5\alpha^2L/(n(\alpha-2)^2)$: at $\alpha=3$ and $L=10$ this is $0.45$ at $n=1000$ and $0.045$ at $n=10^4$; at $\alpha=4$, $L=10$, it is $0.2$ and $0.02$. The Chebyshev bound decays like $1/n$; the Chernoff form would decay exponentially.

\subsection{What small networks do}

\Cref{tab:small} reports the measured Lipschitz constant $K_{n,L}(1)$ at small width. By the scaling identity \cref{eq:scaling}, $K_{n,L}(\alpha)\le1$ exactly when $\alpha\le K_{n,L}(1)^{-2/L}$, so the median of $K_{n,L}(1)^{-2/L}$ is the variance parameter at which half the networks are non-expansive. At $n=2$ the constant is exact, by sweeping all directions of the plane; for $n\ge3$ it is a sampled lower bound, so the threshold column is an upper bound. At depth one the operator norm of $W_1$ approaches $2$ at $\alpha=1$ (measured $1.995$ at $n=2000$), so the threshold approaches $\tfrac14$.

\begin{table}[h]\centering
\begin{tabular}{rrrrrl}
\toprule
$n$ & $L$ & median $K_{n,L}(1)$ & 10\%--90\% & median threshold $K^{-2/L}$ & method\\
\midrule
2 & 1 & 1.253 & 0.72--1.84 & 0.64 & exact\\
2 & 2 & 0.851 & 0.25--1.78 & 1.18 & exact\\
2 & 3 & 0.336 & 0.00--1.52 & 2.07 & exact\\
2 & 4 & 0.065 & 0.00--1.02 & 3.93 & exact\\
3 & 2 & 1.336 & 0.65--2.36 & 0.75 & sampled\\
3 & 3 & 0.876 & 0.24--1.91 & 1.09 & sampled\\
5 & 3 & 1.442 & 0.97--2.25 & 0.78 & sampled\\
8 & 3 & 2.067 & 1.50--2.81 & 0.62 & sampled\\
\bottomrule
\end{tabular}
\caption{The global Lipschitz constant at $\alpha=1$ for small networks (400 draws at $n=2$, fewer at larger $n$).}
\label{tab:small}
\end{table}

Two features of \cref{tab:small} are worth reading against the theorem. At width two, deeper networks die: with probability growing in $L$ some layer maps everything to zero, $K=0$, and the median threshold rises above $2$ for that reason alone, which is a small-width effect the theorem says nothing about. At width eight and depth three the median threshold is about $0.6$, well below $2$, in keeping with the depth $L_0$ the budget demands before the bound takes effect. The theorem is about the joint regime of large width first and large depth after; the table is a reminder of what it does not say.

\subsection{One matrix, many steps}

\Cref{tab:horizon} measures the statistic of \cref{thm:fh}, the worst relative deviation $\max_{t\le8}\abs{2^t\norm{x_t}^2-1}$ over eight steps, for the tied dynamics of \cref{sec:horizon} and, beside it, for the feedforward pass with a fresh matrix at every step (the object of \cref{sec:intro} at $\alpha=1$); 200 draws per width. Both columns shrink with the width, as both theorems say. The tied matrix wanders more than fresh matrices at this horizon: at $N=2000$ six percent of the tied trajectories still stray by more than one half, against half a percent for fresh matrices. The theorem's rate is exponential in $N$ with constants that depend on $T$, and eight steps of one matrix are not eight independent draws.

\begin{table}[h]\centering
\begin{tabular}{rrrrrrr}
\toprule
 & \multicolumn{3}{c}{one matrix} & \multicolumn{3}{c}{fresh matrix per step}\\
$N$ & median & 90\% & frac.\ $>\tfrac12$ & median & 90\% & frac.\ $>\tfrac12$\\
\midrule
20 & 0.983 & 12.93 & 0.935 & 0.917 & 2.663 & 0.965\\
50 & 0.861 & 4.734 & 0.865 & 0.749 & 1.903 & 0.835\\
100 & 0.719 & 2.556 & 0.685 & 0.512 & 1.159 & 0.545\\
500 & 0.348 & 0.789 & 0.270 & 0.274 & 0.547 & 0.125\\
2000 & 0.192 & 0.434 & 0.060 & 0.137 & 0.243 & 0.005\\
\bottomrule
\end{tabular}
\caption{The worst deviation $\max_{t\le8}|2^t\|x_t\|^2-1|$ over eight steps, 200 draws per width, seed fixed.}
\label{tab:horizon}
\end{table}

\Cref{tab:chernoff} evaluates the rate $I(1/\alpha)$ of \cref{cor:chernoff} by a grid search over $s$, and compares the two bounds on $\PP(K_{n,L}(\alpha)\le1)$ at $L=10$, $n=1000$. Near $\alpha=2$ both are useless at this width; from $\alpha=3$ on the exponential bound is smaller by five, fifteen and thirty orders of magnitude.

\begin{table}[h]\centering
\begin{tabular}{rrrr}
\toprule
$\alpha$ & $I(1/\alpha)$ & Chebyshev, $5\alpha^2L/(n(\alpha-2)^2)$ & Chernoff, $Le^{-nI(1/\alpha)}$\\
\midrule
2.2 & 0.00087 & 6.05 & 4.2\\
2.5 & 0.00455 & 1.25 & 0.11\\
3.0 & 0.01395 & 0.45 & $8.7\times10^{-6}$\\
4.0 & 0.03642 & 0.20 & $1.5\times10^{-15}$\\
6.0 & 0.07847 & 0.11 & $8.4\times10^{-34}$\\
\bottomrule
\end{tabular}
\caption{The expansive side at $L=10$, $n=1000$: the Chebyshev bound of \cref{thm:above2} against the exponential bound of \cref{cor:chernoff}.}
\label{tab:chernoff}
\end{table}
